% GENERATED FILE. Edit canonical lessons, then run npm run generate:books.
% canonical-source-hash: fnv1a64:68ea93a531b0360b
% canonical-lessons: BN-C404-aram, BN-C404-ayesh, BN-C404-jolkhabar, BN-C404-jirono, BN-C404-helan-deoya

\chapter{Rest and Comfort}
\label{ch:bn-rest-and-comfort}
\hlchaptermodality{404}

\begin{chapteropening}
\textbf{By the end of this chapter:} \emph{I can say a room is comfortable, take breakfast and rest a while.}

Complete the last lesson of chapter 404: I can say a room is comfortable, take breakfast and rest a while.
\end{chapteropening}

\section[\texorpdfstring{\bn{আরাম} (ārām)}{আরাম (ārām)}]{\bn{আরাম} (ārām) — comfort; ease; rest}

\label{lesson:BN-C404-aram}



\begin{quote}
Before the new one: say the Bengali for a reception desk, then the Bengali for to arrange.
\end{quote}

\subsection*{You'll want to know: \bn{আরাম}}
\textbf{\bn{আরাম}} — \emph{ārām} — \textquotedblleft{}comfort; ease; rest\textquotedblright{}.

\begin{quote}
\bn{লম্বা} \bn{সফরের} \bn{পরে} \bn{বাড়িতে} \bn{আরাম} \bn{লাগল।}
\end{quote}

\begin{quote}
\emph{lômbā sôphôrer pôre bāṛite ārām lāglô.}
\end{quote}

\begin{quote}
After the long trip it felt good to be home.
\end{quote}

\bn{আরাম} \bn{লাগা}, to feel comfortable; \bn{আরামদায়ক}, comfortable, is built on it.

\begin{cousinweb}
\textbf{\bn{আরাম}} (\emph{ārām}) is usually traced to Persian \emph{ārām}, \textquotedblleft{}rest, quiet, ease\textquotedblright{}.
\end{cousinweb}

\begin{grammarlens}[title={What you've built}]
One more word for getting around.
\end{grammarlens}

\subsection*{Your turn}
\begin{itemize}
\raggedright
  \item \emph{Say it:} \emph{ārām}
  \item \emph{Say it:} \emph{ārām}, once more
  \item \emph{Say it:} \emph{lômbā sôphôrer pôre bāṛite ārām lāglô.} — After the long trip it felt good to be home.
  \item \emph{Recall:} say the Bengali for a reception desk, then the Bengali for to arrange, then say \emph{ārām} again
\end{itemize}

\begin{tcolorbox}[breakable,colback=teal!4,colframe=teal!35!black,title={Before you move on}]
What does \bn{আরাম} mean? (\textquotedblleft{}Comfort; ease; rest\textquotedblright{}.) Say it once more.
\end{tcolorbox}
\section[\texorpdfstring{\bn{আয়েশ} (āyesh)}{আয়েশ (āyesh)}]{\bn{আয়েশ} (āyesh) — leisure; ease}

\label{lesson:BN-C404-ayesh}



\begin{quote}
Before the new one: say the Bengali for to arrange, then the Bengali for comfort.
\end{quote}

\subsection*{You'll want to know: \bn{আয়েশ}}
\textbf{\bn{আয়েশ}} — \emph{āyesh} — \textquotedblleft{}leisure; ease\textquotedblright{}.

\begin{quote}
\bn{ছুটির} \bn{দিনে} \bn{অরুণ} \bn{আয়েশ} \bn{করে} \bn{চা} \bn{খায়।}
\end{quote}

\begin{quote}
\emph{chhuṭir dine ôrun āyesh kore chā khāy.}
\end{quote}

\begin{quote}
On holidays Arun drinks his tea at leisure.
\end{quote}

\bn{আয়েশ} \bn{করে}, at leisure, taking one's time.

\begin{cousinweb}
\textbf{\bn{আয়েশ}} (\emph{āyesh}) is usually traced to Arabic \emph{ʿaysh}, \textquotedblleft{}living, a life of ease\textquotedblright{}, which reached Bengali through Persian.
\end{cousinweb}

\begin{grammarlens}[title={What you've built}]
One more word for getting around.
\end{grammarlens}

\subsection*{Your turn}
\begin{itemize}
\raggedright
  \item \emph{Say it:} \emph{āyesh}
  \item \emph{Say it:} \emph{āyesh}, once more
  \item \emph{Say it:} \emph{chhuṭir dine ôrun āyesh kore chā khāy.} — On holidays Arun drinks his tea at leisure.
  \item \emph{Recall:} say the Bengali for to arrange, then the Bengali for comfort, then say \emph{āyesh} again
\end{itemize}

\begin{tcolorbox}[breakable,colback=teal!4,colframe=teal!35!black,title={Before you move on}]
What does \bn{আয়েশ} mean? (\textquotedblleft{}Leisure; ease\textquotedblright{}.) Say it once more.
\end{tcolorbox}
\section[\texorpdfstring{\bn{জলখাবার} (jôlkhābār)}{জলখাবার (jôlkhābār)}]{\bn{জলখাবার} (jôlkhābār) — breakfast; a light meal}

\label{lesson:BN-C404-jolkhabar}



\begin{quote}
Before the new one: say the Bengali for comfort, then the Bengali for leisure.
\end{quote}

\subsection*{You'll want to know: \bn{জলখাবার}}
\textbf{\bn{জলখাবার}} — \emph{jôlkhābār} — \textquotedblleft{}breakfast; a light meal\textquotedblright{}.

\begin{quote}
\bn{হোটেলে} \bn{জলখাবারে} \bn{লুচি} \bn{আর} \bn{আলুর} \bn{তরকারি} \bn{ছিল।}
\end{quote}

\begin{quote}
\emph{hoṭele jôlkhābāre luchi ār ālur tôrkāri chhilô.}
\end{quote}

\begin{quote}
At the hotel there were luchis and potato curry for breakfast.
\end{quote}

Literally water-food: the small meal taken with a glass of water, morning or afternoon.

\begin{cousinweb}
\textbf{\bn{জলখাবার}} (\emph{jôlkhābār}) joins \textbf{\bn{জল}}, \textquotedblleft{}water\textquotedblright{}, Sanskrit \emph{jala} (Monier-Williams), to \textbf{\bn{খাবার}}, \textquotedblleft{}food\textquotedblright{}, from \bn{খাওয়া}, to eat (CDIAL 3865); the whole means a light meal, not food and water.
\end{cousinweb}

\begin{grammarlens}[title={What you've built}]
One more word for getting around.
\end{grammarlens}

\subsection*{Your turn}
\begin{itemize}
\raggedright
  \item \emph{Say it:} \emph{jôlkhābār}
  \item \emph{Say it:} \emph{jôlkhābār}, once more
  \item \emph{Say it:} \emph{hoṭele jôlkhābāre luchi ār ālur tôrkāri chhilô.} — At the hotel there were luchis and potato curry for breakfast.
  \item \emph{Recall:} say the Bengali for comfort, then the Bengali for leisure, then say \emph{jôlkhābār} again
\end{itemize}

\begin{tcolorbox}[breakable,colback=teal!4,colframe=teal!35!black,title={Before you move on}]
What does \bn{জলখাবার} mean? (\textquotedblleft{}Breakfast; a light meal\textquotedblright{}.) Say it once more.
\end{tcolorbox}
\section[\texorpdfstring{\bn{জিরোনো} (jirono)}{জিরোনো (jirono)}]{\bn{জিরোনো} (jirono) — to rest a while; to catch one's breath}

\label{lesson:BN-C404-jirono}



\begin{quote}
Before the new one: say the Bengali for leisure, then the Bengali for breakfast.
\end{quote}

\subsection*{You'll want to know: \bn{জিরোনো}}
\textbf{\bn{জিরোনো}} — \emph{jirono} — \textquotedblleft{}to rest a while; to catch one's breath\textquotedblright{}.

\begin{quote}
\bn{পাহাড়ে} \bn{ওঠার} \bn{পথে} \bn{আমরা} \bn{একটা} \bn{গাছের} \bn{নিচে} \bn{জিরোলাম।}
\end{quote}

\begin{quote}
\emph{pāhāṛe oṭhār pôthe āmrā ekṭā gāchher niche jirolām.}
\end{quote}

\begin{quote}
On the way up the hill we rested under a tree.
\end{quote}

Past (he/she …): \bn{তিনি} \bn{গাছের} \bn{নিচে} \bn{জিরোলেন}, \emph{tini gāchher niche jirolen}.

A short stop on the way; \bn{বিশ্রাম} \bn{নেওয়া} is a longer rest.

\begin{cousinweb}
\textbf{\bn{জিরোনো}} (\emph{jirono}) is inherited: CDIAL lists Bengali \emph{jirāna}, \textquotedblleft{}to rest, to repose\textquotedblright{}, under the reconstructed \emph{jīrayati} (CDIAL 5236).
\end{cousinweb}

\begin{grammarlens}[title={What you've built}]
One more word for getting around.
\end{grammarlens}

\subsection*{Your turn}
\begin{itemize}
\raggedright
  \item \emph{Say it:} \emph{jirono}
  \item \emph{Say it:} \emph{jirono}, once more
  \item \emph{Say it:} \emph{pāhāṛe oṭhār pôthe āmrā ekṭā gāchher niche jirolām.} — On the way up the hill we rested under a tree.
  \item \emph{Recall:} say the Bengali for leisure, then the Bengali for breakfast, then say \emph{jirono} again
\end{itemize}

\begin{tcolorbox}[breakable,colback=teal!4,colframe=teal!35!black,title={Before you move on}]
What does \bn{জিরোনো} mean? (\textquotedblleft{}To rest a while; to catch one's breath\textquotedblright{}.) Say it once more.
\end{tcolorbox}
\section[\texorpdfstring{\bn{হেলান} \bn{দেওয়া} (helān deoyā)}{হেলান দেওয়া (helān deoyā)}]{\bn{হেলান} \bn{দেওয়া} (helān deoyā) — to lean back; to recline}

\label{lesson:BN-C404-helan-deoya}



\begin{quote}
Before the new one: say the Bengali for breakfast, then the Bengali for to rest a while.
\end{quote}

\subsection*{You'll want to know: \bn{হেলান} \bn{দেওয়া}}
\textbf{\bn{হেলান} \bn{দেওয়া}} — \emph{helān deoyā} — \textquotedblleft{}to lean back; to recline\textquotedblright{}.

\begin{quote}
\bn{ট্রেনে} \bn{আমি} \bn{জানালায়} \bn{হেলান} \bn{দিয়ে} \bn{ঘুমালাম।}
\end{quote}

\begin{quote}
\emph{ṭrene āmi jānālāy helān diye ghumālām.}
\end{quote}

\begin{quote}
On the train I leaned against the window and slept.
\end{quote}

Past (he/she …): \bn{তিনি} \bn{চেয়ারে} \bn{হেলান} \bn{দিলেন}, \emph{tini cheyāre helān dilen}.

\begin{cousinweb}
\textbf{\bn{হেলান} \bn{দেওয়া}} (\emph{helān deoyā}) joins \bn{হেলান}, \textquotedblleft{}a leaning\textquotedblright{}, from the verb \bn{হেলা}, which CDIAL lists as Bengali \emph{helā}, \textquotedblleft{}to lean\textquotedblright{}, under Sanskrit \emph{abhilīyatē}, \textquotedblleft{}clings to\textquotedblright{} (CDIAL 529), to \bn{দেওয়া}, to give.
\end{cousinweb}

\begin{grammarlens}[title={What you've built}]
One more word for getting around.
\end{grammarlens}

\subsection*{Your turn}
\begin{itemize}
\raggedright
  \item \emph{Say it:} \emph{helān deoyā}
  \item \emph{Say it:} \emph{helān deoyā}, once more
  \item \emph{Say it:} \emph{ṭrene āmi jānālāy helān diye ghumālām.} — On the train I leaned against the window and slept.
  \item \emph{Recall:} say the Bengali for breakfast, then the Bengali for to rest a while, then say \emph{helān deoyā} again
\end{itemize}

\begin{tcolorbox}[breakable,colback=teal!4,colframe=teal!35!black,title={Before you move on}]
What does \bn{হেলান} \bn{দেওয়া} mean? (\textquotedblleft{}To lean back; to recline\textquotedblright{}.) Say it once more.
\end{tcolorbox}
